% !TeX root = tesis.tex

\subsection{El esquema no interactivo \Groth}
\label{subsec:apendice-b-groth16-esquema-no-interactivo-groth}

En esta sección presentamos una descripción simplificada del argumento de conocimiento subyacente al esquema \Groth, pensada para fijar la intuición algebraica principal sin reproducir todos los detalles técnicos de una formalización estándar. \Groth es un \zkSNARK no interactivo en el modelo de una cadena de referencia común: una vez generado el $\CRS$, el prover produce una única prueba y el verifier la comprueba sin enviarle desafíos. Sea $\Relation \subseteq \InstanceSpace \times \WitnessSpace$ una relación en $\NP$. Como en la Sección~\ref{sec:aritmetizacion-y-construccion-de-un-snark-qap-quadratic-arithmetic-programs}, suponemos que la relación ha sido aritmetizada como un \QAP de tamaño $m$ y dimensión $n$ sobre un cuerpo finito $\K = \Fp$, y todo se reduce a argumentar sucintamente una identidad polinomial. Vamos a ligar todos estos nombres de objetos antes de describir el esquema.

En particular, existen $3n$ polinomios $\{A_i(x)\}_{i=1}^n$, $\{B_i(x)\}_{i=1}^n$, $\{C_i(x)\}_{i=1}^n$ en $\poly{\Fp}_{< m}$ tales que, definiendo el vector extendido de asignación

$$\AssignmentVector = (z_1, \ldots, z_n) \in \Fp^n$$

donde $\AssignmentVector$ representa el vector completo de asignación del \QAP: incluye la coordenada constante, las coordenadas correspondientes a los \textit{public inputs}, las correspondientes al witness original del problema, y las variables auxiliares introducidas por la aritmetización. En particular, adoptamos la convención usual de escribir $x = (x_1, \ldots, x_\ell)$ para la parte pública con $x_1 = 1$, de modo que la constante queda incluida dentro de los \textit{public inputs}; si $w = (w_1, \ldots, w_t)$ denota el witness original, entonces tanto $x$ como $w$ aparecen embebidos dentro de $\AssignmentVector$. Definiendo los polinomios agregados

$$A(x) = \sum_{i=1}^n \AssignmentVector_i A_i(x), B(x) = \sum_{i=1}^n \AssignmentVector_i B_i(x), C(x) = \sum_{i=1}^n \AssignmentVector_i C_i(x),$$

y para un polinomio objetivo $Z$ construido sobre $r_1, \ldots, r_m \in \Fp$ se cumple que $(x, w) \in \Relation$ si y solo si existe un polinomio $H$ tal que

$$AB - C = HZ.$$

Además, sean $\GOne$ y $\GTwo$ grupos cíclicos aditivos de orden primo $p$, sean $G_1 \in \GOne$ y $G_2 \in \GTwo$ generadores, y sea $\GT$ un grupo multiplicativo, con un pairing bilineal no degenerado $e: \GOne \times \GTwo \to \GT$.

\begin{definition}[Esquema no interactivo \Groth]
\label{def:apendice-b-groth16-esquema-no-interactivo-groth}
	El \emph{esquema no interactivo \Groth} busca argumentar la existencia de un polinomio $H$ y consta de los siguientes algoritmos:
	
	\begin{itemize}
		\item $\Setup$: Se eligen de forma uniforme $\tau, \alpha, \beta, \gamma, \delta \in \mathbb{F}^*_p$, que se usan para construir la \SRS, y luego son destruidos. La \SRS contiene:
		\begin{itemize}[noitemsep]
			\item Las potencias de $\tau$ necesarias para evaluar polinomios de hasta grado $m - 1$ en los grupos $\GOne, \GTwo$: $\{\tau^k G_1, \tau^k G_2\}_{k=0}^{m-1}$.
			\item Los elementos $\alpha G_1, \beta G_2, \gamma G_2, \delta G_1, \delta G_2$.
			\item Para $\ell + 1 \leq i \leq n$ (i.e., los índices correspondientes a la parte privada del vector de asignación), los elementos
			$$\frac{\beta A_i(\tau) + \alpha B_i(\tau) + C_i(\tau)}{\delta} G_1.$$
			\item Para los inputs públicos $1 \leq i \leq \ell$, la clave de verificación contiene los elementos
			$$\frac{\beta A_i(\tau) + \alpha B_i(\tau) + C_i(\tau)}{\gamma} G_1,$$
			que permiten al verifier reconstruir la combinación lineal asociada a la instancia pública.
			\item Los elementos $\left\{\frac{\tau^k Z(\tau)}{\delta} G_1\right\}_{k=0}^{m-2}$.
		\end{itemize}
		\item $\Prove$: El prover $\Prover$, que conoce $x$ y $w$, calcula $A(\tau), B(\tau), C(\tau), H(\tau)$. Luego elige aleatoriamente $r, s \in \mathbb{F}^*_p$, y define:
		\begin{enumerate}
			\item $\pi_A = \left(\alpha + A(\tau) + r\delta\right) G_1$.
			\item $\pi_B = \left(\beta + B(\tau) + s\delta\right) G_2$.
			\item $\pi_C = \left(\sum_{i=\ell+1}^{n} \AssignmentVector_i \frac{\beta A_i(\tau) + \alpha B_i(\tau) + C_i(\tau)}{\delta}\right)G_1 + \left(\frac{H(\tau)Z(\tau)}{\delta} + rB(\tau) + sA(\tau) - rs\delta\right)G_1 $.
		\end{enumerate}
			El \textit{output} del algoritmo $\Prove$ es la tripla $\pi = (\pi_A, \pi_B, \pi_C) \in (\GOne, \GTwo, \GOne)$.
			\item $\Verify$: El verifier $\Verifier$, que solamente conoce los inputs públicos $x_1,\ldots,x_\ell$, reconstruye la combinación
			$$V_{\mathrm{pub}}(\tau) := \sum_{i=1}^{\ell} x_i \frac{\beta A_i(\tau) + \alpha B_i(\tau) + C_i(\tau)}{\gamma}.$$
			Luego, acepta si:
			$$e(\pi_A, \pi_B) = e(\alpha G_1, \beta G_2) \cdot e\!\left(V_{\mathrm{pub}}(\tau) G_1, \gamma G_2\right) \cdot e(\pi_C, \delta G_2).$$
		\end{itemize}
	\end{definition}

No está demás recordar que \Groth es un argumento de conocimiento. Intuitivamente, esto significa que para todo prover eficiente que logra producir pruebas aceptadas con probabilidad no despreciable, existe un extractor probabilístico en tiempo polinomial, con acceso a dicho prover, que recupera un witness $w$ tal que $(x, w) \in \Relation$.

El paper original de Groth \cite{cryptoeprint:2016/260} describe el protocolo en términos de vectores y combinaciones lineales, no en términos de exponenciaciones de grupo. La formulación moderna que presentamos (y sobre la cual están basadas las implementaciones principales) interpreta cada elemento del \SRS $a \in \Fp$ como elementos de $\GOne$ mediante la aplicación $a \mapsto aG_1$, de modo que las combinaciones lineales se vuelven \MSM (multi-scalar multiplications) y el verifier pueda usar pairings bilineales.

Observamos que si $(x, w) \in \Relation$, entonces por la definición del \QAP existe un polinomio $H$ tal que $AB - C = HZ$, y en particular se cumple la identidad evaluada en el trapdoor del setup:
$$A(\tau)B(\tau) - C(\tau) = H(\tau)Z(\tau).$$
Si expandimos la ecuación de verificación y utilizamos la bilinealidad de $e$, se observa que todos los términos dependientes de $r$ y $s$ se cancelan, y la igualdad se satisface completamente. Por lo tanto, el protocolo cumple la propiedad de completitud. La solidez se fundamenta en supuestos criptográficos estándar sobre grupos bilineales y en la estructura algebraica de la construcción.

A diferencia de los $\Sigma$-\textit{protocols} discutidos en la Sección~\ref{sec:aritmetizacion-y-construccion-de-un-snark-protocolo-fiatshamir-y-no-interactividad}, \Groth no parte de un intercambio en el que el verifier elige un \textit{challenge}, ni utiliza la transformación de \FiatShamir. El $\Setup$ genera el $\CRS$ a partir de los secretos $(\tau,\alpha,\beta,\gamma,\delta)$ y esos secretos deben destruirse; el $\CRS$ resultante fija la estructura algebraica con la que se producen y verifican las pruebas. Por separado, cada ejecución de $\Prove$ utiliza aleatoriedad fresca $r,s$, responsable de randomizar la prueba y contribuir a la propiedad de conocimiento cero. Algunos análisis de las propiedades algebraicas de \Groth se formulan en el \AGM \cite{Fuchsbauer2018}.

Para finalizar esta subsección, vamos a observar que \Groth realmente está diseñado para verificar la identidad polinomial en el exponente (en $\GT$), pero separando la parte pública $1 \leq i \leq \ell$, de la parte privada $\ell + 1 \leq i \leq n$, y de la parte de aleatorización. En efecto:

$$\begin{aligned}
	\pi_C &= \left(\sum_{i=\ell+1}^{n} \AssignmentVector_i \frac{\beta A_i(\tau) + \alpha B_i(\tau) + C_i(\tau)}{\delta}\right) G_1 \\
	&+ \frac{H(\tau)Z(\tau)}{\delta} G_1 + (rB(\tau) + sA(\tau) - rs\delta)G_1.
\end{aligned}$$

Así, vemos que todo lo que está dividido por $\delta$ corresponde a la parte algebraica del \QAP, mientras que los términos $rB(\tau) + sA(\tau) - rs\delta$ aparecen como una perturbación cuadrática simétrica en $r, s$.


\subsection{\Groth como extensión natural de \KZG}
\label{subsec:apendice-b-groth16-groth-como-extension-natural-de-kzg}

Es instructivo interpretar \Groth como una evolución del esquema \KZG aplicado a la aritmetización mediante \QAP. Recordemos que en \KZG, el prover demuestra que conoce un polinomio $P(x)$ tal que $y = P(z)$. La verificación consiste en chequear, mediante un pairing bilineal, que $P(\tau) - y$ es divisible por $\tau - z$. En esencia, \KZG permite verificar una identidad polinomial en un punto oculto $\tau$.

En el caso del \QAP, lo que queremos verificar es: $A(x) \cdot B(x) - C(x) = H(x) Z(x)$. En vez de argumentar divisibilidad por el polinomio $x - z$, lo hacemos por el polinomio $Z(x)$. Pero \KZG no sería suficiente para esto, por los siguientes tres motivos que \Groth soluciona.

\begin{itemize}[itemsep=0.1cm]
	\item En \KZG, el compromiso revela implícitamente la evaluación $P(\tau)$, sin soportar separación entre parte pública y privada. En cambio, en \Groth el \textit{input} puede ocultarse, y los \textit{public inputs} deben poder verificarse en forma separada. Es por esto que aparecen los parámetros $\alpha, \beta, \gamma, \delta$, que mezclan algebraicamente los polinomios de forma que $\Verifier$ pueda aislar la parte pública, y el \textit{witness} quede enmascarado.
	\item En \KZG, necesitaríamos un compromiso para cada uno de $A(x), B(x), C(x), H(x)$, lo que conlleva varias pruebas de divisibilidad con correspondientes ecuaciones de pairing bilineal. \Groth comprime toda la verificación del \QAP en una única ecuación de emparejamiento.
	\item Como se señaló en la Subsección~\ref{subsec:apendice-b-kzg-lo-que-tenemos-hasta-ahora}, \KZG por sí solo no es \ZK. La aleatoriedad $r, s$ introducida por \Groth, así como los términos cruzados $(rB(\tau) + sA(\tau) - rs\delta)G_1$ garantizan que la distribución de la prueba sea independiente del witness, y por lo tanto es \ZK.
\end{itemize}

\subsection{Complejidad de \Groth}
\label{subsec:apendice-b-groth16-complejidad-de-groth}

En la fase de $\Setup$ (que recordemos, en la práctica es una ceremonia en la que varias personas deben contribuir con aleatoriedad y al menos una debe comportarse de forma honesta), para computar la \SRS en su totalidad hay que computar:
\begin{itemize}[noitemsep]
	\item $\Theta(m)$ multiplicaciones escalares en $\GOne$.
	\item $\Theta(m)$ multiplicaciones escalares en $\GTwo$ (en la práctica, $\GTwo$ suele tener el doble de bits que $\GOne$).
	\item La \SRS tiene una cantidad lineal de elementos en $m$ y $\ell$, con $\ell \leq n$. En particular, a medida que aumenta el tamaño del circuito, aumenta también el tamaño del \textit{trusted setup} necesario. El \textit{trusted setup} de \Groth es una fase costosa, y no es universal, sino específica para cada circuito.
\end{itemize}

En la fase de generación de la prueba, el costo principal proviene de las \MSM necesarias para computar $\pi_A, \pi_B, \pi_C$. Recordemos que una \MSM es una operación del estilo $\sum s_i P_i$, donde los $s_i$ son escalares y los $P_i$ son puntos de curva, de modo que en \Groth las \MSM son aquellas operaciones donde multiplicamos un polinomio evaluado en $\tau$ por el generador correspondiente. Esto permite el uso de algoritmos especializados, principalmente el de Pippenger \cite{Pippenger1976}, usado en las bibliotecas estándar de sistemas \zkSNARK \cite{gnark, arkworks, bellman}. Cada una de ellas cuesta $\Theta(m)$ operaciones escalares en $\GOne$ o $\GTwo$ respectivamente. Por ejemplo:

$$A(\tau) G_1 = \sum_{i=0}^{m - 1} a_i \tau^i G_1.$$

\begin{itemize}[itemsep = 0.1 cm]
	\item Una \MSM en $\GOne$ para calcular $A(\tau) G_1$ durante el cómputo de $\pi_A$.
	\item Una \MSM en $\GTwo$ para calcular $B(\tau) G_2$ durante el cómputo de $\pi_B$.
	\item Una \MSM en $\GOne$ durante el cómputo de $\pi_C$. Acá es donde se observa que los términos de la \SRS correspondientes al witness permiten alivianar el cómputo:
	$$\begin{aligned}
		\pi_C &= \sum_{i=\ell+1}^{n} \AssignmentVector_i \left(\frac{\beta A_i(\tau) + \alpha B_i(\tau) + C_i(\tau)}{\delta} G_1 \right) + \sum_{j=0}^{m-2} h_j \left(\frac{\tau^j Z(\tau)}{\delta} G_1 \right) \\
		&+ \sum_{i=1}^{n} (r\AssignmentVector_i) B_i(\tau) G_1 + \sum_{i=1}^{n} (s\AssignmentVector_i) A_i(\tau) G_1 - rs(\delta G_1).
	\end{aligned},$$
	y más aún, en la práctica se fusionan algebraicamente todas las contribuciones en un único vector de escalares de modo que se hace una \MSM.
	\item En el cálculo de $H(x)$, se requiere una multiplicación polinomial y una división por $Z(x)$. Con una \FFT, esto cuesta $\BigTheta(n \log n)$ operaciones en $\Fp$.
	\item La prueba generada $\pi = (\pi_A, \pi_B, \pi_C)$ tiene tamaño constante $\BigTheta(1)$ (tres elementos de grupo).
\end{itemize}

En la fase de verificación de la prueba, la complejidad temporal del algoritmo no depende del tamaño total del circuito, sino linealmente en el número de public inputs $\ell$:

\begin{itemize}[itemsep=0.1 cm]
	\item Reconstruir la combinación lineal correspondiente a los inputs públicos, lo que es una \MSM en $\GOne$ de tamaño $\ell$, y típicamente $\ell \ll m$.
	\item Evaluar la ecuación de pairing bilineal, para lo cual se requieren evaluar tres pairings (pues $\pairing{\alpha G_1}{\beta G_2}$ es constante desde el $\Setup$) y algunas multiplicaciones en $\GT$.
\end{itemize}

\subsection{Aplicación de \Groth al proyecto}
\label{subsec:apendice-b-groth16-aplicacion-de-groth-al-proyecto}

En contextos de blockchain, la métrica relevante no es únicamente la complejidad asintótica, sino el costo concreto en términos de operaciones criptográficas soportadas por la máquina virtual subyacente. En particular, las operaciones dominantes en la verificación de \Groth son los pairings bilineales. Si la plataforma expone precompilados o primitivas nativas eficientes para emparejamientos (como ocurre en ciertos entornos basados en curvas elípticas tipo \BN o \BLS), el costo de verificación es esencialmente constante y pequeño, lo que hace a \Groth especialmente adecuado para validación \textit{onchain}.

Esta observación es particularmente relevante en el contexto de \Cardano, donde las capacidades criptográficas disponibles \Onchain determinan directamente qué construcciones \zkSNARK usar. La decisión correspondiente se justifica en la Subsección~\ref{subsec:design-decisions-esquema-de-argumentos-zk}, y su integración concreta con los dos circuitos auxiliares se desarrolla en las Secciones~\ref{sec:tx-snapshot-membership-el-circuito-de-pertenencia-a-una-transactionsnapshot} y~\ref{sec:tx-set-update-el-circuito-txsetupdatecircuit}. En particular, las primitivas de pairings bilineales sobre \BLS disponibles en \Plutus, junto con la extrema sucintez del argumento \Groth, determinan que sea el \zkSNARK utilizado en el proyecto \cite{groth_validator_modulop}.
